\chapter{Hidradenitis Suppurativa}
\index{Hidradenitis Suppurativa}
\label{sec:Hidradenitis Suppurativa}

\section{Overview}
Hidradenitis suppurativa (HS) is a chronic, recurrent inflammatory skin condition affecting apocrine gland-bearing skin, primarily the axillae, inguinal region, and perineum. It is characterized by painful nodules, abscesses, draining sinus tracts, and scarring. HS affects approximately 1--4\% of the population worldwide and has a significant impact on quality of life due to pain, drainage, and social stigma.
\section{Classification}

\begin{description}[font=\normalfont\bfseries,leftmargin=3em,labelwidth=2.5em]
  \item[Cause] Inflammatory
  \item[Organ System] Skin
  \item[Pathophysiology] Follicular occlusion with subsequent inflammation, rupture, and chronic suppurative granulomatous response
  \item[Duration] Chronic (years to lifetime)
  \item[Onset] Gradual (post-pubertal)
  \item[Mode of Transmission] Non-communicable
  \item[Clinical Course] Chronic, relapsing-remitting with progressive scarring
  \item[Severity] Mild to Severe
  \item[Extent of Disease] Localized to Regional
  \item[Age Group] Adults (post-pubertal, peak 20--40 years)
  \item[Epidemiology] Prevalence 1--4\%; female predominance 3:1
  \item[ICD-11 Code] ED91.0
\end{description}

\section{Causes}
HS begins with follicular hyperkeratosis and occlusion of the terminal hair follicle, leading to follicular dilation and rupture. Released follicular contents trigger an intense inflammatory response with neutrophils, macrophages, and lymphocytes. Secondary bacterial colonization (often polymicrobial) exacerbates inflammation. Genetic predisposition (autosomal dominant in some families) and dysregulated cytokine signaling (TNF-alpha, IL-17, IL-1 beta) play key roles.

\section{Risk Factors}

\begin{itemize}[noitemsep]
  \item Cigarette smoking (strongly associated)
  \item Obesity and metabolic syndrome
  \item Family history of HS
  \item Female sex
  \item Mechanical friction and skin occlusion
  \item Acne vulgaris (association)
  \item Inflammatory bowel disease (Crohn's disease)
  \item Polycystic ovary syndrome
\end{itemize}

\section{Signs and Symptoms}

\subsection{Early symptoms}
\begin{itemize}[noitemsep]
  \item Early manifestations of hidradenitis suppurativa may be subtle and nonspecific
\end{itemize}

\subsection{Common symptoms}
\begin{itemize}[noitemsep]
  \item \subsection{Early symptoms}
  \item \subsection{Common symptoms}
  \item \textbf{Common symptoms:}
  \item Deep, painful, erythematous nodules in apocrine gland-bearing areas
  \item Recurrent abscesses with purulent drainage
  \item Formation of sinus tracts and tunnels connecting lesions
  \item Hypertrophic scarring and contractures
  \item Chronic malodorous drainage
  \item Pain that is often severe and debilitating
  \item Difficulty performing daily activities
  \item Emotional distress and social isolation
\end{itemize}

\subsection{Less common symptoms}
\begin{itemize}[noitemsep]
  \item Atypical or less common presentations of hidradenitis suppurativa may occur
\end{itemize}

\subsection{Emergency warning signs}
\begin{itemize}[noitemsep]
  \item Severe or worsening symptoms of hidradenitis suppurativa require immediate medical evaluation
\end{itemize}
\section{Progression and Stages}
HS progresses through three Hurley stages. Stage I involves isolated or multiple abscesses without sinus tracts or scarring. Stage II features recurrent abscesses with sinus tracts and scarring, widely separated. Stage III shows diffuse or near-diffuse involvement with interconnected sinus tracts, abscesses, and extensive scarring. Most patients remain at Stage I or II; progression to Stage III is associated with smoking, obesity, and delay in treatment.

\section{Complications}

\begin{itemize}[noitemsep]
  \item Chronic pain and physical disability
  \item Scarring and contractures limiting mobility
  \item Secondary infection and cellulitis
  \item Systemic inflammation and anemia of chronic disease
  \item Squamous cell carcinoma developing in chronic sinus tracts (rare)
  \item Depression, anxiety, and social isolation
  \item Reduced quality of life
  \item Emotional and psychological effects
\end{itemize}

\section{Diagnosis}

\begin{itemize}[noitemsep]
  \item Clinical diagnosis based on characteristic lesion morphology and distribution
  \item Hurley staging to classify severity
  \item Skin biopsy to rule out alternative diagnoses
  \item Bacterial culture of draining lesions for targeted antibiotic therapy
  \item Lab tests to monitor inflammatory markers
\end{itemize}

\section{Prevention}

\begin{itemize}[noitemsep]
  \item Smoking cessation
  \item Weight loss and management of metabolic syndrome
  \item Avoidance of skin friction and occlusive clothing
  \item Maintenance of good skin hygiene
  \item Go for regular check-ups so any problems are caught early
\end{itemize}

\section{Treatment Options}

\begin{itemize}[noitemsep]
  \item Topical clindamycin 1\% solution for mild disease
  \item Oral antibiotics (tetracyclines, clindamycin plus rifampin combination)
  \item Biologic therapy (adalimumab, infliximab) for moderate to severe disease
  \item Hormonal therapy (spironolactone, oral contraceptives)
  \item Intralesional corticosteroids for acute inflammatory nodules
  \item Incision and drainage of acute abscesses
  \item Laser hair removal
  \item Wide surgical excision for refractory localized disease
  \item Lifestyle changes and self-care
\end{itemize}

\section{Living with It}

\begin{itemize}[noitemsep]
  \item Follow your treatment plan as prescribed
  \item Keep up with regular appointments with your healthcare provider
  \item Adjust your daily habits as needed
  \item Reach out to family, friends, or support groups for help
  \item Keep learning about your condition and how to manage it
\end{itemize}

\section{Outlook}
HS is a chronic condition requiring long-term management. With appropriate treatment, many patients achieve significant improvement in symptoms and quality of life. Early diagnosis and treatment improve outcomes and prevent progression to advanced stages. Biologic therapies have dramatically improved outcomes for moderate to severe disease. Without treatment, HS tends to progress and cause permanent scarring.
